\chapter{结论}

VeriSpecOSLab 为生成式 AI 进入操作系统实验课提供了一套完整平台。学生先通过问答理解课程概念和讨论设计，再亲自编写系统、模块、接口、扩展目标和跨模块修改规格。平台完成确定性检查，AI 以只读方式评审现有规格。学生修改并提交后，实现助手根据当前 Lab 和模块范围生成代码与配套测试。平台在独立工作目录中检查实际代码变化，并运行构建和相关测试。QEMU、开发板、报告与提交使用同一份项目配置和规格。

平台已经实现本地 CLI、问答、评审、实现、调试和验证五个 Agent 角色、规格检查、Git 实现流程、结构化 Runner、连续哈希审计、知识库引用、报告生成和提交归档。完整源码的 RISC-V xv6 课程案例覆盖 Lab 1--10，从最小工程、内存、进程和文件系统逐步进入完整 usertests 与 VisionFive 2 四核移植。它验证了平台能够承载一条连续、可运行、可复查的操作系统课程流程。

暑假期间，平台面向华东师范大学计算机拔尖班的部分学生进行了两节课试讲。试讲集中发现三类问题，平台安装与操作还不够方便，学生对课程所需前置知识的了解不够充分，独立进行大型项目开发的经验也不足。我们据此收敛本地学生流程，重写课程入口和 Lab 1，将设计决定分散到对应 Lab，并补充环境检查、分步指导、Agent 实现检查和确定性验证。当前 workspace 类型检查、应用构建与 xv6 Lab 1--10 课程历史审计均通过。

VeriSpecOSLab 将 AI 的代码生成能力与学生主导的理解、设计和验收过程结合起来。平台为学生提供及时帮助，也让课程保留设计决定、代码修改、测试过程和人工复核。暑期试讲已经形成从课堂观察到平台调整的完整反馈闭环，后续课程教学、更多内核案例和更丰富验证方法可以继续沿用同一套规格与证据结构。
